\documentclass[11pt,a4paper]{article}
\usepackage{iftex}
\ifPDFTeX
  \usepackage[utf8]{inputenc}
  \usepackage[T1]{fontenc}
  \usepackage{lmodern}
\else
  \usepackage{fontspec}   % XeLaTeX/LuaLaTeX (e.g. Tectonic): Latin Modern OpenType has all Slovak glyphs
\fi
% shorthands=off: Slovak babel makes "-" active, which breaks \cmidrule{2-6} (booktabs)
\usepackage[slovak,shorthands=off]{babel}
\usepackage[margin=2.5cm]{geometry}
\usepackage{booktabs}
\usepackage{multirow}
\usepackage{graphicx}
\usepackage{amsmath}
\usepackage{xcolor}
\usepackage[hidelinks]{hyperref}
\usepackage[numbers,sort&compress]{natbib}
\bibliographystyle{plainnat}

\newcommand{\todo}[1]{\textcolor{red}{[TODO: #1]}}
\newcommand{\sat}{S@10}
\usepackage{microtype}
\setlength{\emergencystretch}{3em}  % fewer overfull lines with long Slovak words and \texttt names

\title{Viacjazyčné a medzijazykové vyhľadávanie overených tvrdení\\
\large SemEval-2025 Task~7 -- 1.~etapa: návrh projektu a baseline riešenie}
\author{Viktor Hančovský, Benjamín Konkoľ, Jakub Spiśák, Roman Šafranko}
\date{\today}

\begin{document}
\maketitle

\begin{abstract}
Zaoberáme sa úlohou SemEval-2025 Task~7: pre príspevok zo sociálnej siete treba v~kolekcii stoviek tisíc už overených tvrdení (fact-checkov) nájsť tie, ktoré sa ho týkajú, a to aj vtedy, keď sú príspevok a fact-check v~rôznych jazykoch.
V~prvej etape sme implementovali reprodukovateľný experimentálny rámec s~oficiálnym protokolom hodnotenia a dve baseline riešenia: lexikálne BM25 a husté vyhľadávanie pomocou modelu multilingual-E5-large bez dotrénovania.
Na testovacej množine dosahuje E5 baseline \sat{} 0{,}886 (monolingválna úloha) a 0{,}681 (medzijazyková úloha), čím prekonáva oficiálne baseline organizátorov (0{,}859 a 0{,}637).
Upravené BM25 so stemovaním, odstraňovaním stop slov a segmentáciou thajčiny dosahuje 0{,}868 a 0{,}623.
\end{abstract}

% =====================================================================
\section{Úvod}
\label{sec:intro}

\paragraph{Problém.}
Dezinformácie sa šíria rýchlejšie, než ich stihnú profesionálni overovatelia faktov (fact-checkeri) vyvracať.
Veľká časť nepravdivých tvrdení sa však opakuje: to isté tvrdenie sa objavuje v~nových príspevkoch, v~iných formuláciách, na iných platformách a často aj v~iných jazykoch.
Ak by overovateľ vedel pre nový príspevok rýchlo nájsť už existujúci fact-check, nemusí tvrdenie overovať znova.
Takáto úloha sa nazýva \emph{vyhľadávanie už overených tvrdení} (previously fact-checked claim retrieval) \citep{shaar-etal-2020-known}.

\paragraph{Formulácia úlohy.}
Pre príspevok $p$ zo sociálnej siete (text, prípadne text extrahovaný z~priložených obrázkov pomocou OCR) a kolekciu fact-checkov $\mathcal{F}$ má systém vrátiť usporiadaný zoznam fact-checkov tak, aby sa relevantné fact-checky nachádzali čo najvyššie.
Ide o~úlohu \emph{vyhľadávania a usporiadania} (ranking), nie o~klasifikáciu pravdivosti.
Úloha SemEval-2025 Task~7 \citep{peng-etal-2025-semeval} má dve podúlohy:
\begin{itemize}
  \item \textbf{monolingválnu} --- príspevok aj fact-check sú v~rovnakom jazyku a kolekcia kandidátov je obmedzená na fact-checky v~jazyku príspevku,
  \item \textbf{medzijazykovú} (crosslingual) --- príspevok a fact-check môžu byť v~rôznych jazykoch a kolekcia kandidátov nie je obmedzená.
\end{itemize}

\paragraph{Motivácia.}
Fact-checky existujú prevažne vo veľkých jazykoch (najmä v~angličtine), zatiaľ čo dezinformácie sa šíria vo všetkých jazykoch.
Medzijazykové vyhľadávanie umožňuje znovu využiť existujúcu prácu overovateľov aj v~jazykoch, kde je ich málo \citep{pikuliak-etal-2023-multilingual, peng-etal-2025-semeval}.

\paragraph{Ciele projektu.}
\begin{enumerate}
  \item Vytvoriť reprodukovateľný experimentálny rámec s~oficiálnym protokolom hodnotenia (1.~etapa).
  \item Implementovať a vyhodnotiť lexikálne a husté baseline riešenie (1.~etapa).
  \item Navrhnúť vlastné riešenie, ktoré baseline podstatne zlepší: fúziu výsledkov, dotrénovanie modelu na trénovacích pároch so ťažkými negatívnymi príkladmi a preusporiadanie kandidátov (2.~etapa).
  \item Systematicky analyzovať chyby a výsledky interpretovať po jazykoch (2.~etapa).
\end{enumerate}
Projekt sa zameriava primárne na monolingválnu podúlohu; medzijazykovú podúlohu vyhodnocujeme paralelne, keďže zdieľa celý systém aj predvypočítané reprezentácie.

% =====================================================================
\section{Prehľad existujúcich prístupov}
\label{sec:related}

\paragraph{Vyhľadávanie už overených tvrdení.}
\citet{shaar-etal-2020-known} úlohu formulovali a použili dvojfázový prístup: kandidátov získali pomocou BM25 a preusporiadali ich modelom založeným na BERT s~učením sa usporiadania (learning to rank).
\citet{vo-lee-2020-facts} vyhľadávali fact-checkové články k~tweetom s~využitím textu aj obrázkov, čo je relevantné aj pre nás, keďže časť príspevkov obsahuje text iba v~obrázku.
\citet{shaar-etal-2022-role} ukázali, že kontext na strane vstupného tvrdenia v~politickej debate môže zlepšiť vyhľadávanie o~viac ako 10 bodov MAP; nejde o~priamy experiment s~OCR príspevkov zo sociálnych sietí.
\citet{kazemi-etal-2021-claim} rozšírili párovanie tvrdení na jazyky mimo angličtiny (hindčina, bengálčina, malajálamčina, tamilčina) pomocou destilovaného viacjazyčného modelu viet.
\citet{pikuliak-etal-2023-multilingual} vytvorili dátovú množinu MultiClaim (približne 28\,tis. príspevkov v~27 jazykoch a 206\,tis. fact-checkov v~39 jazykoch), z~ktorej vychádza aj táto úloha.
Najlepším modelom bez dotrénovania bol v~ich porovnaní GTR-T5-Large \citep{ni-etal-2022-large}.

\paragraph{Reprezentácie textu pre vyhľadávanie.}
Klasickým lexikálnym modelom je BM25 \citep{robertson-zaragoza-2009-bm25}.
Husté vyhľadávanie používa bi-enkodéry, ktoré zakódujú dotaz a dokument nezávisle do vektora a podobnosť počítajú ako kosínusovú podobnosť \citep{reimers-gurevych-2019-sentence}.
Viacjazyčné bi-enkodéry vznikajú destiláciou z~jednojazyčných modelov \citep{reimers-gurevych-2020-making} alebo kontrastívnym predtrénovaním na veľkých viacjazyčných korpusoch, napríklad modely E5 \citep{wang-etal-2022-e5, wang-etal-2024-multilingual-e5} a BGE-M3 \citep{chen-etal-2024-m3}.
Presnejšie, ale pomalšie sú cross-enkodéry, ktoré spracujú dotaz aj dokument naraz a používajú sa na preusporiadanie malého počtu kandidátov \citep{nogueira-cho-2019-passage}.
Výsledky viacerých vyhľadávačov sa dajú kombinovať pomocou reciprocal rank fusion \citep{cormack-etal-2009-rrf}.

\paragraph{Riešenia v~SemEval-2025 Task~7.}
Súťaže sa zúčastnilo 179 registrovaných účastníkov s~52 testovacími odovzdaniami \citep{peng-etal-2025-semeval}.
Poradia a nasledujúce monolingválne skóre preberáme z~oficiálneho prehľadu (tabuľka~8); články tímov môžu uvádzať iné zaokrúhlenie.
Najlepší výsledok bol \sat{} 0{,}9601, desiate miesto 0{,}9232.
Zdokumentované systémy často kombinujú husté vyhľadávanie s~preusporiadaním alebo fúziou modelov, existujú však aj silné systémy bez týchto krokov.
\begin{itemize}
  \item \textbf{TIFIN India} \citep{devadiga-etal-2025-tifin} (3.~miesto, 0{,}9383): preklad príspevkov do angličtiny modelom Aya-expanse-8B, dotrénovanie Stella-400M s~ťažkými negatívnymi príkladmi, preusporiadanie 50 kandidátov modelom Qwen2.5-72B a RRF.
  Samotný dotrénovaný vyhľadávač s~prekladmi dosiahol 0{,}9217 (tabuľka~5 článku).
  \item \textbf{RACAI} \citep{chivereanu-tufis-2025-racai} (4.~miesto, 0{,}9377): BGE-Multilingual-Gemma2-9B na originálnych textoch, kontrastívne dotrénovaný pomocou LoRA, bez preusporiadania.
  Prompt tuning, QLoRA a Matryoshka reprezentácie skúmali samostatne; nemožno ich všetky pripísať finálnej konfigurácii.
  \item \textbf{QUST\_NLP} \citep{liu-etal-2025-qust} (5.~miesto, 0{,}9365): výber vyhľadávačov podľa jazyka, preusporiadanie a vážené hlasovanie.
  V~abláciách modelov E5 spojenie originálu s~prekladom zlepšilo monolingválny priemer, zatiaľ čo v~medzijazykovej úlohe bol lepší samotný preklad (\S4.1, tabuľky~1--2 článku); nejde o~pravidlo pre každý model a jazyk.
  \item \textbf{UniBuc-AE} \citep{enache-2025-unibuc} (6.~miesto, 0{,}9336): anglický e5-large-v2 a viacjazyčný multilingual-e5-large-instruct, kontrastívne učenie a následné dotrénovanie s~ťažkými negatívnymi príkladmi.
  Štyri konfigurácie kombinujú váženými skóre podobnosti, s~váhami vybranými na validačných dátach.
  \item \textbf{UWBa} \citep{lenc-etal-2025-uwba} (7.~miesto, 0{,}9270): bez dotrénovania, nad anglickými prekladmi.
  Finálny monolingválny systém podľa jazyka vyberá NV-Embed-v2 samostatne alebo v~dvojici s~GPT či Mistral; medzijazykový systém používa iba NV-Embed.
  Päť modelov skúmali pri predbežnom výbere, nie vo finálnej spoločnej kombinácii.
  \item \textbf{fact check AI} \citep{rastogi-2025-fact} (10.~miesto, 0{,}9232): dotrénované bi-enkodéry a kombinovanie modelov trénovaných v~piatich foldoch.
  Monolingválne spracovanie využíva originály aj preklady, medzijazykové anglické preklady.
  \item \textbf{Duluth} \citep{syed-pedersen-2025-duluth}: finálny lexikálny systém TF-IDF so slovnou tokenizáciou a~15\,000 príznakmi dosiahol 0{,}6883 (tabuľka~3 článku; zaokrúhlene 0{,}69).
\end{itemize}
Tieto výsledky motivujú samostatne overiť prínos dotrénovania, preusporiadania a fúzie na našej vývojovej množine; nezaručujú zlepšenie pre každú konfiguráciu.
Výpočtové nároky publikovaných veľkých modelov nemožno preniesť na našu 8\,GB GPU bez samostatného overenia; v~2.~etape preto plánujeme menšie modely a dotrénovanie na vzdialenej GPU.

% =====================================================================
\section{Dáta}
\label{sec:data}

Používame oficiálnu dátovú množinu úlohy \citep{semeval2025task7-dataset}, ktorá je dostupná na výskumné účely na základe žiadosti.
Trénovacia a vývojová časť vychádza z~MultiClaim \citep{pikuliak-etal-2023-multilingual}; testovacia časť obsahuje nové dáta zozbierané rovnakou metodikou.

\paragraph{Štruktúra.}
Každý \emph{príspevok} obsahuje text napísaný používateľom a texty extrahované z~priložených obrázkov (OCR, Google Vision API).
Každý \emph{fact-check} obsahuje overované tvrdenie (claim) a názov článku (title).
Všetky texty sú dostupné v~origináli aj v~anglickom preklade (Google Translate) spolu s~detegovaným jazykom.
Rozdelenie na podúlohy, trénovacie, vývojové a testovacie príspevky aj kolekciu kandidátov pre každú podúlohu definuje súbor \texttt{tasks.json}.
Monolingválna úloha má v~trénovacej a vývojovej časti 8 jazykov (ara, deu, eng, fra, msa, por, spa, tha); testovacia časť pridáva poľštinu a turečtinu, ktoré neboli vopred ohlásené.

% Source: results/data_stats.md, generated by scripts/data_stats.py
\begin{table}[t]
\centering
\caption{Základné štatistiky monolingválnej úlohy a v~spodnej časti medzijazykovej úlohy: počty príspevkov, veľkosť kolekcie kandidátov (podľa \texttt{tasks.json}), počet správnych párov a podiel hodnotených príspevkov s~viac ako jedným správnym fact-checkom.}
\label{tab:data}
\setlength{\tabcolsep}{4pt}
\begin{tabular}{l rrrrr @{\hspace{12pt}} rrrr}
\toprule
 & \multicolumn{5}{c}{Trénovacia a vývojová časť} & \multicolumn{4}{c}{Testovacia časť} \\
\cmidrule(lr){2-6}\cmidrule(lr){7-10}
Jazyk & train & dev & kolekcia & páry & $>$1 (\%) & test & kolekcia & páry & $>$1 (\%) \\
\midrule
ara & 676 & 78 & 14\,201 & 78 & 0{,}0 & 500 & 21\,153 & 514 & 2{,}8 \\
deu & 667 & 83 & 4\,996 & 101 & 21{,}7 & 500 & 7\,485 & 574 & 10{,}4 \\
eng & 4\,351 & 478 & 85\,734 & 627 & 21{,}1 & 500 & 145\,287 & 740 & 26{,}4 \\
fra & 1\,596 & 188 & 4\,355 & 200 & 4{,}8 & 500 & 6\,316 & 549 & 8{,}4 \\
msa & 1\,062 & 105 & 8\,424 & 116 & 9{,}5 & 93 & 686 & 96 & 3{,}2 \\
pol & -- & -- & -- & -- & -- & 500 & 8\,796 & 564 & 9{,}4 \\
por & 2\,571 & 302 & 21\,569 & 403 & 28{,}1 & 500 & 32\,598 & 613 & 18{,}8 \\
spa & 5\,628 & 615 & 14\,082 & 692 & 11{,}4 & 500 & 25\,440 & 576 & 12{,}8 \\
tha & 465 & 42 & 382 & 42 & 0{,}0 & 183 & 583 & 183 & 0{,}0 \\
tur & -- & -- & -- & -- & -- & 500 & 12\,536 & 550 & 9{,}0 \\
\midrule
spolu & 17\,016 & 1\,891 & 153\,743 & 2\,259 & 15{,}5 & 4\,276 & 260\,880 & 4\,959 & 11{,}5 \\
\midrule
\multicolumn{10}{l}{\emph{Medzijazyková úloha}} \\
všetky & 4\,972 & 552 & 153\,743 & 651 & 14{,}3 & 4\,000 & 272\,447 & 5\,322 & 20{,}5 \\
\bottomrule
\end{tabular}
\end{table}

\begin{table}[t]
\centering
\caption{Medzijazyková úloha: najčastejšie kombinácie jazyka príspevku a jazyka správneho fact-checku (podiel zo všetkých správnych párov). Jazyk je detegovaný (Google Translate), nie anotovaný.}
\label{tab:crosslang}
\begin{tabular}{lr @{\hspace{18pt}} lr}
\toprule
\multicolumn{2}{c}{Vývojová časť (651 párov)} & \multicolumn{2}{c}{Testovacia časť (5\,322 párov)} \\
\cmidrule(r{12pt}){1-2}\cmidrule(l){3-4}
príspevok $\to$ fact-check & \% & príspevok $\to$ fact-check & \% \\
\midrule
hin $\to$ eng & 15{,}2 & hin $\to$ eng & 43{,}5 \\
spa $\to$ spa & 14{,}0 & eng $\to$ eng & 20{,}2 \\
eng $\to$ eng & 13{,}4 & ara $\to$ ara & 4{,}4 \\
por $\to$ por & 8{,}9 & spa $\to$ spa & 4{,}2 \\
zho $\to$ eng & 7{,}4 & eng $\to$ hin & 3{,}2 \\
sin $\to$ eng & 6{,}9 & ben $\to$ eng & 2{,}7 \\
urd $\to$ eng & 5{,}5 & tur $\to$ tur & 2{,}7 \\
tgl $\to$ eng & 5{,}1 & urd $\to$ eng & 2{,}6 \\
\midrule
fact-checky v~angličtine & 60{,}7 & fact-checky v~angličtine & 73{,}9 \\
\bottomrule
\end{tabular}
\end{table}

\begin{table}[t]
\centering
\caption{Zloženie a dĺžka textov (vývojová a testovacia časť, hodnotené príspevky oboch podúloh). Dĺžka je počet slov oddelených medzerami v~anglickom preklade (medián / 90.~percentil).}
\label{tab:composition}
\begin{tabular}{lcc}
\toprule
 & Vývojová časť & Testovacia časť \\
\midrule
Hodnotených príspevkov & 2\,443 & 8\,276 \\
Iba text používateľa (\%) & 45{,}6 & 53{,}3 \\
Iba text z~obrázka -- OCR (\%) & 9{,}8 & 0{,}0 \\
Text aj OCR (\%) & 44{,}6 & 46{,}7 \\
Dĺžka príspevku, monolingválna & 45 / 198 & 45 / 187 \\
Dĺžka príspevku, medzijazyková & 42 / 198 & 39 / 141 \\
Dĺžka fact-checku (tvrdenie + názov) & 26 / 41 & 27 / 46 \\
\bottomrule
\end{tabular}
\end{table}


\paragraph{Pozorovania.}
\begin{itemize}
  \item \textbf{Veľkosti kolekcií sa líšia o~tri rády}: od 382 (thajčina) po 145\,287 (angličtina, test).
  Úloha je preto v~rôznych jazykoch rôzne ťažká a priemerovanie cez jazyky dáva malým jazykom rovnakú váhu ako angličtine.
  Niektoré vývojové množiny sú veľmi malé (thajčina 42, arabčina 78 príspevkov), takže rozdiely v~jednotlivých jazykoch na vývojovej množine treba interpretovať opatrne.
  \item \textbf{Približne 11--21\,\% príspevkov má viac správnych fact-checkov} (rovnaké tvrdenie overilo viacero organizácií), čo metriku \sat{} mierne zvýhodňuje.
  \item \textbf{Text v~obrázkoch je dôležitý}: takmer polovica príspevkov obsahuje OCR a vo vývojovej časti má 9{,}8\,\% príspevkov text \emph{iba} v~obrázku.
  V~testovacej časti takéto príspevky nie sú, preto vplyv OCR na vývojovej množine môže byť väčší ako na testovacej.
  \item \textbf{Medzijazyková testovacia množina je nevyvážená}: podľa detegovaného jazyka tvoria hindské príspevky 40{,}9\,\% a páry hindčina$\to$angličtina 43{,}5\,\% všetkých správnych párov; 73{,}9\,\% správnych fact-checkov je v~angličtine.
  Vo vývojovej časti je podiel hindčiny len 14{,}3\,\%, takže výsledky na vývojovej a testovacej množine nemusia byť priamo porovnateľné (tab.~\ref{tab:crosslang}).
  Zaujímavé je, že značná časť párov v~medzijazykovej úlohe má podľa detektora \emph{rovnaký} jazyk príspevku aj fact-checku (napr. eng$\to$eng 20\,\% na teste); medzijazyková úloha sa teda od monolingválnej líši najmä neobmedzenou kolekciou kandidátov.
  \item \textbf{Fact-checky sú krátke} (medián 26 slov), príspevky dlhšie a variabilnejšie (90.~percentil približne 190 slov), čo pri modeloch s~limitom 512 tokenov vedie k~orezaniu časti príspevkov.
\end{itemize}

% =====================================================================
\section{Metodika}
\label{sec:method}

\paragraph{Metriky.}
Hlavnou metrikou úlohy je \emph{Success@10} (\sat{}): podiel príspevkov, pre ktoré je medzi prvými desiatimi vrátenými fact-checkmi aspoň jeden správny.
Metrika zodpovedá praktickému použitiu, keď overovateľ prezrie krátky zoznam kandidátov.
Doplnkovo uvádzame \emph{Mean Reciprocal Rank} orezaný na 10 (MRR@10), ktorý zohľadňuje aj pozíciu prvého správneho výsledku.
V~monolingválnej úlohe sa \sat{} počíta pre každý jazyk zvlášť a výsledok je ich priemerom (každý jazyk má rovnakú váhu bez ohľadu na počet príspevkov); v~medzijazykovej úlohe sa počíta cez všetky príspevky naraz.
Takto počítajú výsledky aj oficiálne hodnotiace skripty.

\paragraph{Protokol hodnotenia.}
\begin{itemize}
  \item Všetky návrhové rozhodnutia (variant textu, nastavenia analyzátora) robíme \textbf{iba na vývojovej množine}.
  Na testovacej množine spúšťame každú vybranú konfiguráciu práve raz a výsledky slúžia len na finálne porovnanie.
  \item Kolekciu kandidátov pre každý príspevok berieme presne podľa \texttt{tasks.json}: v~monolingválnej úlohe fact-checky jazyka príspevku, v~medzijazykovej všetkých 153\,743 (vývoj) resp. 272\,447 (test) fact-checkov.
  \item Naša implementácia metrík je overená proti oficiálnym hodnotiacim skriptom: na štyroch oficiálnych baseline predikciách organizátorov dáva presne rovnaké hodnoty ako uvádza článok organizátorov.
  Každý súbor s~predikciami vyhodnocujeme oboma implementáciami.
  \item Všetky experimenty sú riadené konfiguračnými súbormi, majú fixované náhodné semienka a ukladajú metriky aj použitú konfiguráciu.
\end{itemize}

\paragraph{Reprezentácia textu.}
Text príspevku tvorí používateľský text spojený so všetkými OCR textami; text fact-checku tvorí tvrdenie spojené s~názvom.
Porovnávame tri varianty: originálny text (\emph{orig}), anglický preklad (\emph{eng}) a ich spojenie (\emph{both}).

% =====================================================================
\section{Baseline riešenia}
\label{sec:baseline}

\paragraph{BM25.}
Implementovali sme BM25 (Okapi, $k_1 = 1{,}5$, $b = 0{,}75$) pomocou riedkych matíc, aby bolo vyhľadávanie v~kolekcii s~270\,tis. dokumentmi rýchle (5--25\,s na celý beh).
Výsledky sú overené testom zhody s~knižnicou \texttt{rank\_bm25}.
Základná verzia používa tokenizáciu na slová (\texttt{\textbackslash w+}) bez ďalšieho spracovania.
Upravená verzia pridáva jazykovo špecifické stemovanie (Snowball), odstraňovanie stop slov a segmentáciu thajčiny pomocou knižnice PyThaiNLP \citep{phatthiyaphaibun-etal-2023-pythainlp}.
Thajčina nepoužíva medzery medzi slovami a samohláskové a tónové znamienka nepatria do triedy \texttt{\textbackslash w}, takže základná tokenizácia rozbije thajský text na nezmyselné fragmenty.

\paragraph{Husté vyhľadávanie (multilingual-E5-large).}
Používame model \texttt{intfloat/\allowbreak multilingual-\allowbreak e5-large} \citep{wang-etal-2024-multilingual-e5} bez dotrénovania, s~predponami \texttt{query:} a \texttt{passage:} podľa odporúčania autorov a maximálnou dĺžkou vstupu 512 tokenov.
Reprezentácie všetkých fact-checkov vypočítame raz (fp16, GPU RTX 5070 s~8\,GB) a uložíme na disk; výpočet pre 270\,tis. fact-checkov trvá približne 5--7 minút.
Vyhľadávanie je skalárny súčin normalizovaných vektorov (menej ako 1\,s na celú množinu).

% =====================================================================
\section{Výsledky}
\label{sec:results}

% Source: results/report_tables.md (Table A), generated by scripts/make_report_tables.py
\begin{table}[t]
\centering
\caption{Hlavné výsledky na testovacej množine. Monolingválna \sat{} a MRR@10 sú priemerom cez 10 jazykov, medzijazykové sú počítané cez všetky príspevky. Hodnoty organizátorov sme prepočítali z~ich zverejnených predikcií (\sat{} sa zhoduje s~publikovanými hodnotami na $10^{-4}$). Najlepší výsledok v~stĺpci je zvýraznený.}
\label{tab:main}
\begin{tabular}{lcccc}
\toprule
 & \multicolumn{2}{c}{Monolingválna} & \multicolumn{2}{c}{Medzijazyková} \\
\cmidrule(lr){2-3}\cmidrule(lr){4-5}
Systém & \sat{} & MRR@10 & \sat{} & MRR@10 \\
\midrule
\multicolumn{5}{l}{\emph{Baseline organizátorov}} \\
BM25 & 0{,}812 & 0{,}653 & 0{,}603 & 0{,}382 \\
paraphrase-multilingual-mpnet-base-v2 & 0{,}568 & 0{,}395 & 0{,}411 & 0{,}253 \\
GTR-T5-large & 0{,}730 & 0{,}503 & \textbf{0{,}705} & \textbf{0{,}442} \\
multilingual-E5-large & 0{,}859 & 0{,}688 & 0{,}637 & 0{,}395 \\
\midrule
\multicolumn{5}{l}{\emph{Naše baseline}} \\
BM25, základná tokenizácia (eng) & 0{,}771 & 0{,}622 & 0{,}537 & 0{,}339 \\
BM25 + stem. + stop slová + thajčina\textsuperscript{a} & 0{,}868 & 0{,}713 & 0{,}623 & 0{,}382 \\
multilingual-E5-large\textsuperscript{b} & \textbf{0{,}886} & \textbf{0{,}725} & 0{,}681 & 0{,}417 \\
\bottomrule
\end{tabular}

\smallskip
{\footnotesize \textsuperscript{a}~Monolingválna: variant \emph{both}; medzijazyková: \emph{eng} so stemovaním a stop slovami. \textsuperscript{b}~Monolingválna: \emph{both}; medzijazyková: \emph{eng}. Varianty boli vybrané na vývojovej množine.}
\end{table}


% Source: results/report_tables.md (Table C)
\begin{table}[t]
\centering
\caption{Ablácie na vývojovej množine (\sat{}). Na základe týchto výsledkov sme vybrali konfigurácie pre testovaciu množinu (zvýraznené). Varianty textu: \emph{orig} = originál, \emph{eng} = anglický preklad, \emph{both} = spojenie.}
\label{tab:dev}
\setlength{\tabcolsep}{3.5pt}
\resizebox{\textwidth}{!}{%
\begin{tabular}{llcccccccccc}
\toprule
 & & \multicolumn{9}{c}{Monolingválna} & Medzijaz. \\
\cmidrule(lr){3-11}\cmidrule(lr){12-12}
Model & Konfigurácia & ara & deu & eng & fra & msa & por & spa & tha & priem. & \sat{} \\
\midrule
BM25 & eng & 0{,}705 & 0{,}518 & 0{,}546 & 0{,}750 & 0{,}619 & 0{,}735 & 0{,}724 & 0{,}881 & 0{,}685 & 0{,}413 \\
BM25 & orig & 0{,}731 & 0{,}530 & 0{,}540 & 0{,}761 & 0{,}752 & 0{,}785 & 0{,}766 & 0{,}738 & 0{,}700 & -- \\
BM25 & eng + stem & 0{,}731 & 0{,}578 & 0{,}569 & 0{,}771 & 0{,}724 & 0{,}752 & 0{,}753 & 0{,}905 & 0{,}723 & -- \\
BM25 & eng + stem + stop & 0{,}782 & 0{,}747 & 0{,}630 & 0{,}793 & 0{,}733 & 0{,}811 & 0{,}807 & 0{,}881 & 0{,}773 & \textbf{0{,}520} \\
BM25 & orig + thai & 0{,}731 & 0{,}530 & 0{,}540 & 0{,}761 & 0{,}752 & 0{,}785 & 0{,}766 & 0{,}857 & 0{,}715 & -- \\
BM25 & orig + stem + stop + thai & 0{,}782 & 0{,}723 & 0{,}626 & 0{,}824 & 0{,}733 & 0{,}775 & 0{,}846 & 0{,}905 & 0{,}777 & -- \\
BM25 & both + stem + stop + thai & 0{,}795 & 0{,}759 & 0{,}626 & 0{,}819 & 0{,}790 & 0{,}808 & 0{,}842 & 0{,}905 & \textbf{0{,}793} & 0{,}489 \\
\midrule
mE5 & eng & 0{,}808 & 0{,}759 & 0{,}797 & 0{,}862 & 0{,}790 & 0{,}772 & 0{,}854 & 0{,}976 & 0{,}827 & \textbf{0{,}591} \\
mE5 & orig & 0{,}808 & 0{,}783 & 0{,}780 & 0{,}846 & 0{,}790 & 0{,}798 & 0{,}807 & 0{,}952 & 0{,}821 & 0{,}534 \\
mE5 & both & 0{,}859 & 0{,}819 & 0{,}795 & 0{,}867 & 0{,}895 & 0{,}838 & 0{,}868 & 0{,}952 & \textbf{0{,}862} & 0{,}533 \\
\midrule
\multicolumn{2}{l}{\# príspevkov (dev)} & 78 & 83 & 478 & 188 & 105 & 302 & 615 & 42 & 1\,891 & 552 \\
\bottomrule
\end{tabular}}
\end{table}


% Source: results/report_tables.md (Table B)
\begin{table}[t]
\centering
\caption{\sat{} po jazykoch na testovacej množine (monolingválna úloha). Poľština (pol) a turečtina (tur) sa v~trénovacích dátach nenachádzajú.}
\label{tab:lang}
\setlength{\tabcolsep}{4pt}
\resizebox{\textwidth}{!}{%
\begin{tabular}{lccccccccccc}
\toprule
Systém & ara & deu & eng & fra & msa & pol & por & spa & tha & tur & priemer \\
\midrule
BM25 (org.) & 0{,}846 & 0{,}812 & 0{,}750 & 0{,}822 & 0{,}914 & 0{,}732 & 0{,}754 & 0{,}796 & 0{,}907 & 0{,}790 & 0{,}812 \\
mE5 (org.) & 0{,}916 & \textbf{0{,}872} & \textbf{0{,}818} & 0{,}854 & 0{,}860 & 0{,}798 & 0{,}808 & 0{,}856 & 0{,}929 & 0{,}878 & 0{,}859 \\
BM25 upravené (naše) & 0{,}930 & 0{,}860 & 0{,}760 & \textbf{0{,}898} & 0{,}957 & 0{,}814 & \textbf{0{,}816} & \textbf{0{,}862} & 0{,}945 & 0{,}840 & 0{,}868 \\
mE5 (naše) & \textbf{0{,}942} & \textbf{0{,}872} & 0{,}814 & \textbf{0{,}898} & \textbf{0{,}978} & \textbf{0{,}830} & 0{,}802 & \textbf{0{,}862} & \textbf{0{,}984} & \textbf{0{,}880} & \textbf{0{,}886} \\
\midrule
\# príspevkov & 500 & 500 & 500 & 500 & 93 & 500 & 500 & 500 & 183 & 500 & 4\,276 \\
veľkosť kolekcie & 21\,153 & 7\,485 & 145\,287 & 6\,316 & 686 & 8\,796 & 32\,598 & 25\,440 & 583 & 12\,536 & \\
\bottomrule
\end{tabular}}
\end{table}


\paragraph{Interpretácia.}
\begin{itemize}
  \item \textbf{Husté vyhľadávanie je lepšie ako lexikálne}, najmä v~medzijazykovej úlohe (0{,}681 oproti 0{,}623 na teste), kde majú príspevok a fact-check často málo spoločných slov.
  \item \textbf{Spracovanie textu je pre BM25 zásadné.}
  Odstránenie stop slov bolo najväčším jednotlivým zlepšením (napr. nemčina na vývojovej množine z~0{,}578 na 0{,}747) a segmentácia thajčiny zlepšila thajčinu z~0{,}738 na 0{,}857 pri originálnom texte.
  Upravené BM25 (0{,}868) prekonáva oficiálne BM25 organizátorov (0{,}812) aj ich E5 baseline (0{,}859).
  \item \textbf{Variant textu závisí od podúlohy.}
  V~monolingválnej úlohe pomáha spojenie originálu s~prekladom (E5: 0{,}862 oproti 0{,}827 len s~angličtinou na vývojovej množine), zatiaľ čo v~medzijazykovej úlohe škodí (0{,}533 oproti 0{,}591).
  Rovnaký smer rozdielov uvádza tím QUST\_NLP vo svojich abláciách modelov E5 \citep{liu-etal-2025-qust}.
  Pravdepodobné vysvetlenie (zatiaľ neoverené): v~monolingválnej úlohe originál zachytí presné formulácie a preklad pridá robustnosť; v~medzijazykovej úlohe sú fact-checky prevažne v~inom jazyku ako príspevok, preto originálny text do reprezentácie pridáva skôr šum.
  \item \textbf{Angličtina je najťažší jazyk} (E5: 0{,}814 na teste) napriek tomu, že ide o~najväčší jazyk s~najkvalitnejšími dátami.
  Pravdepodobnou príčinou je veľkosť kolekcie: 145\,tis. kandidátov oproti stovkám až desaťtisícom v~ostatných jazykoch, takže medzi kandidátmi je oveľa viac podobných fact-checkov (overíme v~analýze chýb).
  \item \textbf{Naše E5 prekonáva E5 organizátorov} (0{,}886 oproti 0{,}859).
  Zisk je sústredený v~jazykoch, kde pomáha spojenie originálu a prekladu (malajčina +0{,}118, thajčina +0{,}055, francúzština +0{,}044); v~angličtine, portugalčine a nemčine sú výsledky rovnaké.
  \item \textbf{Obmedzenie:} pri variante \emph{both} sa 11\,\% príspevkov nezmestí do 512 tokenov a ich koniec sa oreže.
\end{itemize}

% =====================================================================
\section{Plán 2.~etapy}
\label{sec:plan}

\begin{enumerate}
  \item \textbf{Fúzia} BM25 a hustého vyhľadávania pomocou RRF \citep{cormack-etal-2009-rrf}: lexikálne a husté vyhľadávanie robia odlišné chyby.
  \item \textbf{Dotrénovanie bi-enkodéra} (mE5) kontrastívnou stratou na trénovacích pároch (približne 20\,tis. monolingválnych a 6\,tis. medzijazykových) so ťažkými negatívnymi príkladmi získanými z~našich baseline výsledkov; podľa literatúry ide o~najväčší zdroj zlepšenia \citep{devadiga-etal-2025-tifin, enache-2025-unibuc}.
  Trénovanie prebehne na GPU T4 (Google Colab).
  \item \textbf{Preusporiadanie} 50 najlepších kandidátov cross-enkodérom (napr. \texttt{bge-\allowbreak reranker-\allowbreak v2-m3}).
  \item \textbf{Vlastný príspevok:} riešenie orezávania dlhých príspevkov a využitie metadát, najmä časovej blízkosti príspevku a fact-checku, ako ďalšieho signálu pri usporiadaní.
  \item \textbf{Analýza chýb} s~vopred definovanou taxonómiou, nezávislou anotáciou dvoma členmi tímu a overením navrhnutej úpravy na oddelených dátach.
  \item \textbf{Demo} pre externých používateľov, ktoré zobrazuje skóre a odkaz na verejný fact-check (nie texty z~dátovej množiny).
\end{enumerate}
Každý krok vyhodnotíme na vývojovej množine ako ablačný experiment voči predchádzajúcemu najlepšiemu systému.

% =====================================================================
\section{Reprodukovateľnosť}
\label{sec:repro}
Kód je v~repozitári \url{https://github.com/Mekr0t/FCCR}.
Súbor \texttt{README.md} popisuje inštaláciu, umiestnenie dát a spúšťanie experimentov.
Každý experiment je definovaný konfiguračným súborom v~\texttt{configs/} a výsledky sú v~\texttt{results/}.
Testovacia sada (pytest) overuje načítanie dát, metriky, zhodu BM25 s~referenčnou knižnicou a správanie analyzátora.
Dátová množina nie je súčasťou repozitára kvôli licencii a treba si ju vyžiadať od organizátorov.

\bibliography{../../docs/references}

\appendix
\section{Použitie LLM a AI agentov}
\label{app:llm}
\todo{Preložiť a vložiť obsah súboru \texttt{docs/llm\_usage.md} (nástroj, model, účel, prompty, výstupy).}

\end{document}
